\section{A Polynomial-Time Sampler for Dense Bridges}
\label{sec:rank-recursive-sampler}

This section gives the main positive result.  It does not analyze a Markov
chain on identity words.  Instead, it recursively separates transpositions
that avoid the largest label from star transpositions incident to that label.
Approximation enters only when long endpoint walks are replaced by their
parity-uniform limits and when a negligible set of recursively short branches
is discarded.  Every sampled word still has exactly the prescribed length and
endpoint.

For $g\in S_r$, write
\[
  \Omega_r(L;g)
  =
  \{(t_1,\ldots,t_L)\in T_r^L:t_1\cdots t_L=g\},
  \qquad
  U_{r,L,g}=\operatorname{Unif}\Omega_r(L;g).
\]
The fiber is nonempty exactly when
\[
  \operatorname{sgn}(g)=(-1)^L
  \quad\hbox{and}\quad
  L\ge r-c(g),
  \tag{RR.1}
\]
where $c(g)$ counts all cycles, including fixed points.  Necessity follows
from sign and cycle count.  Sufficiency follows by writing each cycle as a
minimal product of transpositions and padding by adjacent pairs.

\begin{theorem}[Dense rank-recursive sampler]
\label{thm:dense-rank-recursion}
For every sufficiently large $N$ and every $g\in S_N$ for which
$\Omega_N(2N^3;g)$ is nonempty, there is a randomized algorithm
$\mathsf{DenseSample}(N,2N^3,g)$ with the following properties.
\begin{enumerate}
\item Every output consists of exactly $2N^3$ transpositions and has product
  exactly $g$.
\item If $\nu_{N,g}$ is its output law, then
  \[
    d_{\TV}\bigl(\nu_{N,g},U_{N,2N^3,g}\bigr)
    \le 3N\exp(-N/12288).
    \tag{RR.2}
  \]
\item Its expected Turing runtime and space are polynomial in $N$ and the
  length.  A deliberately conservative bound is
  $O(N^4L^3)$ exact-arithmetic operations on polynomial-bit integers and
  rationals, with $L=2N^3$.
\end{enumerate}
In particular, taking $g=\id$ gives a polynomial-time statistically uniform
sampler for dense identity bridges.
\end{theorem}

The proof is organized around four ingredients: an exact rank
disintegration, exact star completion, a pointwise-smoothed one-rank lift, and
a retained-length telescope.

\subsection{Exact rank disintegration}

Embed $H=S_{r-1}$ in $S_r$ as the stabilizer of $r$.  Call a transposition
\emph{old} if it lies in $H$, and call
\[
  \tau_i=(i\,r),\qquad 1\le i<r,
\]
a \emph{star} transposition.  In the group algebra put
\[
  C_r=\sum_{t\in T_r}t,\qquad
  C_{r-1}=\sum_{t\in T_{r-1}}t,\qquad
  J_r=\sum_{i<r}\tau_i.
\]

\begin{lemma}[Stable-sort bijection]
\label{lem:stable-sort}
Fix an old/star pattern of length $L$ containing $k$ star positions.  There is
a product-preserving bijection between words with this pattern and pairs
$(u,v)$, where $u$ is an old word of length $L-k$ and $v$ is a star word of
length $k$.  If $h=\operatorname{eval}(u)$ and the original endpoint is $g$,
then
\[
  \operatorname{eval}(v)=h^{-1}g.
  \tag{RR.3}
\]
The bijection and its inverse use at most $k(L-k)$ adjacent rewrites.
\end{lemma}

\begin{proof}
For $a\in H$,
\[
  \tau_i a=a\tau_{a^{-1}(i)}.
  \tag{RR.4}
\]
Repeatedly use (RR.4) to move every old factor left of every star factor.
Each rewrite reduces the number of star--old inversions by one, so the
procedure terminates in at most $k(L-k)$ steps and preserves the product.
It is stable: no two old ancestors and no two star ancestors are exchanged.

The inverse starts from $uv$ and restores the chosen binary pattern.  Whenever
a star must move left across an old factor, use the inverse equality
\[
  a\tau_j=\tau_{a(j)}a.
\]
The two procedures undo one another rewrite by rewrite.  Finally
$g=\operatorname{eval}(u)\operatorname{eval}(v)=h\operatorname{eval}(v)$,
which proves (RR.3).
\end{proof}

Since $C_r=C_{r-1}+J_r$ and $C_{r-1}$ commutes with $J_r$, Lemma
\ref{lem:stable-sort} gives the positive coefficient identity
\[
  [g]C_r^L
  =
  \sum_{k=0}^L\binom{L}{k}
  \sum_{h\in H}
    [h]C_{r-1}^{L-k}\,[h^{-1}g]J_r^k.
  \tag{RR.5}
\]
Every summand counts actual words.  Thus an exact sampler could draw $(k,h)$
proportionally to the summand, sample the two endpoint fibers, choose a
uniform binary pattern, and apply the inverse stable sort.  The difficulty is
to draw the labeled branch $(k,h)$ without enumerating $(r-1)!$ possibilities.

\subsection{Exact prescribed-endpoint star bridges}

Write
\[
  B_{r,k}(q)=[q]J_r^k.
\]
The following proposition turns these coefficients into an exact conditional
sampler.  Its enumerative input is the positive transitive
star-factorization formula of F\'eray \cite[Proposition 1.1, Theorem 1.2, and
Lemma 3.1]{Feray2012}.

\begin{proposition}[Star completion]
\label{prop:star-completion}
Given $r,k$ and a labeled endpoint $q\in S_r$, the number $B_{r,k}(q)$ can be
computed in time polynomial in $r$ and $k$.  Whenever it is positive, one can
sample exactly and uniformly from the corresponding length-$k$ star words in
polynomial expected time.
\end{proposition}

\begin{proof}
Let the cycle of $q$ containing $r$ have length $a$, let
$b_1,\ldots,b_m$ be the lengths of the other nontrivial cycles, and let $f$
be the number of fixed points in $\{1,\ldots,r-1\}$.  Put
\[
  \mu=(a,b_1,\ldots,b_m),\quad
  n_0=a+\sum_jb_j,\quad
  \ell_*(q)=n_0+\ell(\mu)-2,
\]
and
\[
  F(x)=\frac{2\sinh(x/2)}{x}
      =\sum_{\substack{d\ge0\\d\ {\rm even}}}
        \frac{x^d}{2^d(d+1)!}.
\]
For $0\le z\le f$, set
\[
  \mu_z=\mu\cup 1^z,\qquad n_z=n_0+z,\qquad
  d_z=k-\ell_*(q)-2z.
\]
For $n_z\ge2$, define
\[
  T_k(\mu_z)
  =
  \frac{k!}{n_z!}
  \left(\prod_{u\in\mu_z}u\right)
  [x^{d_z}]F(x)^{n_z-2}\prod_{u\in\mu_z}F(ux),
  \tag{RR.6}
\]
with a zero contribution when $d_z<0$.  Use the singleton convention
$T_0((1))=1$ and $T_k((1))=0$ for $k>0$.  Decomposing a star word by the
subset of fixed lower labels that it actually uses gives
\[
  B_{r,k}(q)=\sum_{z=0}^f\binom{f}{z}T_k(\mu_z).
  \tag{RR.7}
\]
Indeed, on its exact used-label set the star factor graph is connected, so
F\'eray's transitive formula gives (RR.6); conversely every transitive
factorization on that set uses exactly those labels.  This is a disjoint,
positive decomposition.

All series may be truncated at degree $k$.  There are $O(r)$ factors and
$O(r)$ support choices, so schoolbook polynomial multiplication evaluates
(RR.7) with polynomially many operations on polynomial-bit rationals.

For sampling, partition by the last displayed factor:
\[
  B_{r,j}(q)=\sum_{i<r}B_{r,j-1}(q\tau_i).
  \tag{RR.8}
\]
Choose $i$ with weight $B_{r,j-1}(q\tau_i)$, recurse on
$(j-1,q\tau_i)$, and append $\tau_i$.  Induction using (RR.8) shows that
every word has probability $1/B_{r,j}(q)$.  Exact integer range draws make
the procedure Las Vegas.  A conservative implementation uses
$O(r^3k^3)$ exact-rational operations, and all intermediate integers and
rationals have polynomial bit length.
\end{proof}

\subsection{A positive one-rank lift}

For a rank-$r$ call of length $L$, put $s=L-k$,
\[
  d=\binom{r-1}{2},\qquad b=r-1,
\]
and define
\[
  A_s(h)=[h]C_{r-1}^s,\qquad B_k(q)=[q]J_r^k.
\]
The exact branch weight from (RR.5) is
\[
  \binom{L}{k}A_s(h)B_k(h^{-1}g).
  \tag{RR.9}
\]

We need pointwise, parity-aware estimates rather than ordinary total-variation
mixing.  The following form will be used twice.

\begin{lemma}[Pointwise parity smoothing]
\label{lem:pointwise-smoothing}
Let $K$ be the normalized random-transposition operator associated with the
edges of a connected graph, acting on the corresponding symmetric group.
Suppose its nontrivial absolute spectral radius, after removing the constant
and sign eigenvectors, is at most $\rho$.  On the parity class reachable at
time $t$,
\[
  \left|
    \frac{K^t(e,x)}{2/|G|}-1
  \right|
  \le \frac{|G|-2}{2}\rho^t.
  \tag{RR.10}
\]
For the old and star walks at rank $r$, respectively,
\[
  \rho_{\rm old}=\frac{r-4}{r-2},
  \qquad
  \rho_{\rm star}=\frac{r-2}{r-1}.
  \tag{RR.11}
\]
\end{lemma}

\begin{proof}
Project the two point masses in the matrix coefficient $K^t(e,x)$ away from
the constant and sign eigenspaces.  Cauchy--Schwarz and the operator norm on
the remaining space bound the absolute remainder by
$(1-2/|G|)\rho^t$.  Dividing by the parity-uniform mass $2/|G|$ gives
(RR.10).

For an edge-transposition walk, the interchange-process spectral gap equals
the one-particle graph gap by Caputo--Liggett--Richthammer
\cite[Theorem 1.1]{CLR2010}.  The unnormalized one-particle gaps of
$K_{r-1}$ and the star $K_{1,r-1}$ are $r-1$ and $1$; division by their
numbers of edges gives (RR.11).  Multiplication by the sign character
anticommutes with every transposition operator, so positive and negative
eigenvalues are paired and the same bounds hold in absolute value.
\end{proof}

Consequently, on reachable endpoints,
\begin{align}
  \left|
  \frac{A_s(h)}{2d^s/(r-1)!}-1
  \right|
  &\le
  \delta_{\rm old}(r,s)
  :=
  \frac{(r-1)!-2}{2}
  \left(\frac{r-4}{r-2}\right)^s,
  \tag{RR.12}\\
  \left|
  \frac{B_k(q)}{2b^k/r!}-1
  \right|
  &\le
  \delta_{\rm star}(r,k)
  :=
  \frac{r!-2}{2}
  \left(\frac{r-2}{r-1}\right)^k.
  \tag{RR.13}
\end{align}

The positive lift flattens only the longer side of (RR.9).  Let
$x=g^{-1}(r)$, and let $Z_r(k,x)$ be the number of length-$k$ star words with
endpoint $q$ satisfying $q(x)=r$.  These are exactly the endpoints for which
$h=gq^{-1}$ fixes $r$.  The count has the two-state recurrence
\[
  u_0=1,\quad v_0=0,\qquad
  u_{j+1}=bv_j,\quad v_{j+1}=u_j+(b-1)v_j,
  \tag{RR.14}
\]
with $Z_r(k,r)=u_k$ and $Z_r(k,x)=v_k$ for $x\ne r$.  The same completion
counts give an exact backward sampler for a uniform star word satisfying the
token condition.

After multiplying by one harmless common factor, use the integer $k$-weights
\[
  \widetilde W_k=
  \begin{cases}
  r\binom{L}{k}d^sZ_r(k,x),&s\ge k,\\[3pt]
  \binom{L}{k}d^sb^k,&k>s.
  \end{cases}
  \tag{RR.15}
\]
If $s\ge k$, sample the star token bridge, let its endpoint be $q$, set
$h=gq^{-1}$, and recursively request an old bridge with endpoint $h$.  If
$k>s$, sample the old word unconditionally, let its endpoint be $h$, and use
Proposition \ref{prop:star-completion} at endpoint $h^{-1}g$; this branch does
not recurse.  In both cases choose a uniform binary pattern and invert the
stable sort.

\begin{lemma}[One-rank comparison]
\label{lem:one-rank-comparison}
Let
\[
  \Delta(r,L)=
  \max\left\{
  \delta_{\rm old}\!\left(r,\left\lceil L/2\right\rceil\right),
  \delta_{\rm star}\!\left(r,\left\lfloor L/2\right\rfloor+1\right)
  \right\}.
  \tag{RR.16}
\]
If $\Delta(r,L)<1$, then the exact and approximate branch supports agree, and
the one-rank lift with an exact lower child differs from $U_{r,L,g}$ by at
most
\[
  \frac{\Delta(r,L)}{1-\Delta(r,L)}
  \tag{RR.17}
\]
in total variation, uniformly over feasible $g$.
\end{lemma}

\begin{proof}
In the old-long half, (RR.12) compares every exact old endpoint count with
the flat value used in (RR.15), while the star-coset weight remains exact.
In the star-long half, (RR.13) compares every exact star endpoint count with
the flat value, while the old-word weight remains exact.  Thus every exact
unnormalized branch weight divided by its approximate counterpart lies in
$[1-\Delta,1+\Delta]$.  The same interval contains the ratio of normalizers.
The standard likelihood-ratio estimate then gives
\[
  d_{\TV}(P,Q)
  \le
  \frac{1}{2(1-\Delta)}
  \mathbb E_Q\left|
    \frac{w}{\widetilde w}
    -\frac{Z}{\widetilde Z}
  \right|
  \le \frac{\Delta}{1-\Delta}.
\]
Applying identical exact conditional kernels for old words, star words,
patterns, and inverse sorting cannot increase total variation.
\end{proof}

\subsection{Retained lengths and the rank telescope}

One-rank smoothing is not enough: a recursive old child could be too short.
The following explicit floors retain enough old length at every rank while
discarding only negligible bridge mass.

Fix the top rank $N\ge6$ and set
\[
  a_N=\lceil\log_2N\rceil,\qquad
  z_N=\left\lceil\sqrt{Na_N}\right\rceil.
\]
For $6\le j\le N$, define
\[
  w_j=\frac{1}{j\lceil\sqrt j\rceil}+\frac{100z_N}{Nj},
  \qquad
  W_N=\sum_{j=6}^Nw_j,
  \tag{RR.18}
\]
\[
  c_5=1,\qquad
  c_r=1+\frac{\sum_{j=6}^r w_j}{W_N},
  \qquad
  B_4=0,\qquad
  B_r=\left\lceil Nr(r-1)c_r\right\rceil.
  \tag{RR.19}
\]
In particular $c_N=2$ and
\[
  B_N=2N^2(N-1)<2N^3.
  \tag{RR.20}
\]
At rank $r\ge6$, retain every star-long branch and retain an old-long branch
only if its old length $s$ is at least $B_{r-1}$.  Rank five needs no
truncation, and ranks at most four have exact elementary samplers.

\begin{lemma}[Uniform retained-rank bounds]
\label{lem:retained-rank}
Suppose $N$ is large enough that
\[
  \frac{100z_N(1+\log N)}{N}\le\frac{15}{8},
  \qquad 12z_N\le N,\qquad \log N\le\frac N4.
  \tag{RR.21}
\]
At rank five, for every $L\ge B_5$, there is no discarded branch,
$\Delta(5,L)\le\exp(-5N/4)$, and every recursive child is handled exactly at
rank four.  Moreover, for every $6\le r\le N$, every $L\ge B_r$, and every feasible
$g\in S_r$:
\begin{enumerate}
\item under $U_{r,L,g}$, the mass of discarded old-long branches is at most
  $\exp(-N/12288)$;
\item $\Delta(r,L)\le\exp(-Nr/4)$;
\item every retained recursive child lies in the next retained domain.
\end{enumerate}
\end{lemma}

\begin{proof}
The rank-five statement follows directly from (RR.12)--(RR.16) with
$B_5=20N$: the two nontrivial radii are $1/3$ and $3/4$, and rank-four
children are exact.  We therefore turn to $r\ge6$.

The decreasing-series and harmonic bounds, together with (RR.21), give
$W_N\le3$.  In an unconditioned rank-$r$ word, the old count $A$ has law
\[
  A\sim\operatorname{Bin}\left(L,\frac{r-2}{r}\right).
\]
For a fixed cutoff, its lower-tail probability decreases with $L$, so take
$L=B_r$.  Equations (RR.18)--(RR.19), including the ceiling loss, give
\[
  \mathbb E A-B_{r-1}
  \ge
  \frac{N\sqrt r}{32}+6z_Nr=:t_0.
  \tag{RR.22}
\]
Equivalently, the star count has mean at most $5Nr$.  Bernstein's inequality,
(RR.21), and $(x+y)^2\ge x^2+y^2$ yield
\[
  \Pr[A<B_{r-1}]
  \le
  \exp\!\left(-\frac{N}{12288}-3r\log N\right).
  \tag{RR.23}
\]

This unconditioned estimate must still be transferred to the bridge.  Let
$M_r=\binom r2$.  Lemma \ref{lem:pointwise-smoothing}, applied to the full
transposition walk, gives on the feasible parity class
\[
  P_r^L(e,g)
  \ge
  \frac{2}{r!}
  -
  \left(1-\frac{2}{r-1}\right)^L
  \ge \frac1{r!}
  \tag{RR.24}
\]
for $L\ge B_r$ under (RR.21).  Therefore
$|\Omega_r(L;g)|\ge M_r^L/r!$.  The number of discarded words ending at
$g$ is at most the total number of unconditioned words satisfying the tail
event.  Dividing by (RR.24) and using $r!\le N^r$ in (RR.23) proves
\[
  U_{r,L,g}(\text{discarded})
  \le
  \exp\!\left(-\frac{N}{12288}-2r\log N\right)
  \le \exp(-N/12288).
  \tag{RR.25}
\]

For the second assertion, substitute $L\ge B_r$ into (RR.12)--(RR.16), use
$1-x\le e^{-x}$, and absorb the factorial prefactors using
$\log r\le N/4$.  This gives $\Delta(r,L)\le e^{-Nr/4}$.  The third assertion
is immediate from the retention rule $s\ge B_{r-1}$; star-long branches do
not recurse.
\end{proof}

\begin{proof}[Proof of Theorem \ref{thm:dense-rank-recursion}]
At a retained call $(r,L,g)$, compute the integer weights (RR.15), setting to
zero precisely the discarded old-long weights.  Draw $k$ exactly, execute the
corresponding old-long or star-long procedure, choose a uniform binary
pattern, and invert Lemma \ref{lem:stable-sort}.  The endpoint is always
exactly $g$.

Let $D_r$ be the worst total-variation error over retained rank-$r$ calls.
Ranks at most four are exact: $S_2$ has one feasible word; the old parity
classes at ranks three and four are singletons or exactly uniform, and the
remaining all-star case uses Proposition \ref{prop:star-completion}.  Hence
$D_4=0$.  Comparing first with the exact bridge conditioned on retention and
then using Lemma \ref{lem:one-rank-comparison}, Lemma
\ref{lem:retained-rank} gives
\[
  D_r
  \le
  \exp(-N/12288)+2\exp(-Nr/4)+D_{r-1}.
  \tag{RR.26}
\]
The last term follows from contraction of total variation under the exact
reconstruction kernel.  Summing (RR.26) from rank five to rank $N$ gives
\[
  D_N
  \le
  \sum_{r=5}^N
  \left(\exp(-N/12288)+2\exp(-Nr/4)\right)
  \le 3N\exp(-N/12288).
\]
Equation (RR.20) places the top-level length $2N^3$ in the retained domain.

There are $O(L)$ branch weights, each with $O(L\log N)$ bits.  Token
sampling, permutation updates, and stable sorting are polynomial.  Along one
recursive path there is at most one star-long branch, hence at most one call
to the $O(r^3L^3)$ star sampler.  Summing over ranks gives the stated
conservative $O(N^4L^3)$ exact-arithmetic bound on polynomial-bit operands.
Exact fair-bit range rejection accounts for the expected-runtime qualifier.
\end{proof}

\subsection{Complete transposition support}

The dense target itself contains every transposition with overwhelming
probability.  This is a bridge-conditioned statement, not an ordinary coupon
collector estimate.

\begin{theorem}[Dense complete support]
\label{thm:dense-complete-support}
Let $M=\binom N2$, let $N\ge3$, and let $L$ be even.  Under the uniform
identity bridge,
\[
  \Pr[\text{some transposition is absent}]
  \le
  M\left(1-\frac1M\right)^L
  +
  \frac{MN!}{2}\left(1-\frac2N\right)^L.
  \tag{RR.27}
\]
At $L=2N^3$ this is at most $\exp(-2N)$.  The same conclusion is immediate
for $N=2$.
\end{theorem}

\begin{proof}
Fix a transposition $a$.  Let $P$ be the normalized full-transposition
operator and $Q_a$ the normalized operator with $a$ deleted.  If
$p_L=P^L(e,e)$ and $q_L=Q_a^L(e,e)$, exact word counting gives
\[
  \Pr[a\text{ absent}]
  =
  \left(\frac{M-1}{M}\right)^L\frac{q_L}{p_L}.
  \tag{RR.28}
\]
The graph $K_N-a$ has one-particle Laplacian gap $N-2$.  By
\cite[Theorem 1.1]{CLR2010}, sign symmetry, and the trace formula for the
even return probability,
\[
  q_L\le\frac2{N!}
  +
  \left(1-\frac{N-2}{M-1}\right)^L,
  \qquad
  p_L\ge\frac2{N!}.
\]
Insert these estimates into (RR.28), use
\[
  \frac{M-1}{M}
  \left(1-\frac{N-2}{M-1}\right)
  =1-\frac2N,
\]
and take a union bound over the $M$ choices of $a$.  This proves (RR.27).
At $L=2N^3$, the two terms are each at most $e^{-3N}$ by
$1-x\le e^{-x}$, $\log M\le N$, and $\log(N!)\le N^2$.
\end{proof}

\begin{corollary}
At length $2N^3$, the sampler of Theorem
\ref{thm:dense-rank-recursion} fails to contain every transposition with
probability at most
\[
  \exp(-2N)+3N\exp(-N/12288).
\]
\end{corollary}

\begin{proof}
Apply Theorem \ref{thm:dense-complete-support} to the target and use the
defining event bound for total variation.
\end{proof}

\subsection{Scope and proof status}

Theorem \ref{thm:dense-rank-recursion} is not an exact sampler and not an
arbitrary-accuracy scheme at fixed $(N,L)$.  Its error is a negligible
function of the scheduled parameters, but additional runtime does not lower
the bound below (RR.2).  The remaining standalone problem is to sample the
positive labeled branch law (RR.9), exactly or to arbitrary requested
$\epsilon$, in time polynomial in $r,L,$ and $\log(1/\epsilon)$ without
enumerating $S_{r-1}$ or partition-scale endpoint types.

The result is also polynomial in the permutation degree $N$ and word length
$L$.  If $S_N$ is used to model a reversible function on $n=\log_2N$ wires,
this is not automatically polynomial in $n$.

This proof is a streamlined consolidation of a multi-agent derivation.  Its
algebraic identities and finite boundary cases passed independent exact
implementations, and the retained-length estimates were separately audited.
The section is nevertheless newly assembled and should receive ordinary
human referee review, with particular attention to (RR.22)--(RR.26), before
being treated as a final external theorem.
